% Round 8 author contribution: replace the old lem:sliver and thm:T1 blocks.
% Requires amsmath, amssymb, amsthm and the existing theorem/lemma environments.
% Keep the repaired thm:lemAdp and the analytic thm:T2 in the target document.
% Their directly necessary proofs are included in proof.md, Appendices A and B.
% No old floating-point certificate or T3 numerical enclosure is used here.

\subsection{A whole-region analytic bound for the deep sliver}

\begin{lemma}[Pole-free phase-speed gap bound]\label{lem:r8-phase-speed}
For every $R\ge1$ and $0<u<1/2$, put
$\varepsilon=R^{-1/2}$, $\ell=1/2-u$, and $w=u/\varepsilon$.
Then
\begin{equation}
 G(R,u)\ge
 \frac{\pi^2}{2\varepsilon(w+\ell)(w+\varepsilon\ell)}.
 \label{eq:r8-global-gap-bound}
\end{equation}
No restriction on the second heavy-layer phase is imposed in this statement.
\end{lemma}

\begin{proof}
Let $k=\sqrt\lambda>0$ and normalize the left Dirichlet solution by
$y(0)=0$, $y'(0)=1$. At the interface its vector is
$(y'(u),ky(u))=(\cos(wk),\varepsilon\sin(wk))$.
Propagation across the light interval gives
\begin{align*}
 y'(1/2)&=\cos(wk)\cos(\ell k)
          -\varepsilon\sin(wk)\sin(\ell k),\\
 ky(1/2)&=\varepsilon\sin(wk)\cos(\ell k)
          +\cos(wk)\sin(\ell k).
\end{align*}
Define $A_\varepsilon$ to be the continuous real-valued argument of
$\cos\theta+i\varepsilon\sin\theta$, starting at
$A_\varepsilon(0)=0$. This vector is nonzero, and
\[
 A_\varepsilon'(\theta)=
 \frac{\varepsilon}{\cos^2\theta+\varepsilon^2\sin^2\theta},
 \qquad A_\varepsilon(\pi/2)=\pi/2,
 \qquad A_\varepsilon(\pi)=\pi.
\]
The half-interval phase is
\[
 \Phi(k)=\ell k+A_\varepsilon(wk),\qquad
 0<\Phi'(k)\le\ell+w/\varepsilon.
\]
Its first Neumann and first Dirichlet roots satisfy, respectively,
\[
 \Phi(k_1)=\pi/2,\qquad \Phi(k_2)=\pi.
\]
They are the square roots of the first two full-interval eigenvalues.
Indeed, the spatial phase increases in both layers; these two
half-interval eigenfunctions are positive in their interiors.
Even reflection of the Neumann solution has no interior zero, and
odd reflection of the Dirichlet solution has exactly the midpoint zero.
Sturm oscillation therefore identifies their full-interval indices as
one and two.

Integration of the phase derivative between these actual roots yields
\[
 k_2-k_1\ge\frac{\pi}{2(\ell+w/\varepsilon)}.
\]
Since $\Phi(\pi/(2w))>\pi/2$, one has $wk_1<\pi/2$.
On this first branch,
$A_\varepsilon(\theta)=\arctan(\varepsilon\tan\theta)\le\theta$,
because $\varepsilon\le1$. Hence
\[
 k_1\ge\frac{\pi}{2(\ell+w)}.
\]
Multiplying these bounds gives
\[
 G=\varepsilon^{-2}(k_2-k_1)(k_2+k_1)
 \ge\frac{\pi^2}{2\varepsilon(w+\ell)(w+\varepsilon\ell)}.
\]
Every division above has a positive denominator. In particular, this
argument includes $w<1/2$ and does not assume $\theta_2>\pi/2$.
\end{proof}

\begin{lemma}[Deep sliver, analytic replacement]\label{lem:sliver}
For every $R\ge1500$ and every $0<u\le2/\sqrt R$,
\begin{equation}
 G(R,u)>\frac{15\pi^2}{4}>25,
 \qquad G(R,u)>3\pi^2.
 \label{eq:r8-sliver-margin}
\end{equation}
\end{lemma}

\begin{proof}
Here $w\le2$, $\ell<1/2$ and $\varepsilon<1/38$, since $38^2<1500$.
Equation~\eqref{eq:r8-global-gap-bound} therefore implies
\[
 G>\frac{\pi^2}{2(1/38)(5/2)(2+1/76)}
   =\frac{2888}{765}\pi^2>\frac{15}{4}\pi^2.
\]
The last rational comparison is $11552>11475$.
Also $\pi>3$ gives $15\pi^2/4>135/4>25$, and $15/4>3$.
This bound covers the entire continuous region, including arbitrarily
small positive $w$, $w=1/2$, and the infinite $R$ tail.
\end{proof}

\subsection{T1: the infimum limit and convergence of near-minimizers}

\begin{theorem}[T1, analytic replacement]\label{thm:T1}
Let $u^*$ be the unique minimizer in Theorem~\ref{thm:T2}, and put
$M=\bar D(u^*)$. Then
\[
 \lim_{R\to\infty} Rm_R=M.
\]
More generally, let $R_j\to\infty$, $0<u_j<1/2$, and $\eta_j\ge0$
satisfy
\[
 D_{R_j}(u_j)\le m_{R_j}+\eta_j,
 \qquad R_j\eta_j\longrightarrow0.
\]
Then $u_j\to u^*$. This includes $\eta_j=R_j^{-2}$.
Moreover, $0\le Rm_R-M=O(R^{-1})$.
\end{theorem}

\begin{proof}
The analytic structure in Theorem~\ref{thm:T2} gives
\[
 \bar D(u)\longrightarrow+\infty\quad(u\downarrow0),
 \qquad
 \bar D(u)\longrightarrow3\pi^2\quad(u\uparrow1/2).
\]
Strict increase to the right of $u^*$ implies $M<3\pi^2$.
This inequality uses no numerical enclosure from T3.

First fix any $u\in(0,1/2)$ and set $b=\ell/u$.
For sufficiently large $R$ one has $R\ge1500$ and $w\ge2$.
Use the repaired large-$w$ branch justification and
Theorem~\ref{thm:lemAdp}, whose scope is precisely this region.
Write $\theta_i=u\sqrt{R\lambda_i}$ and $a=a(u)$.
On this branch, $0<\theta_1<\pi/2<\theta_2<\pi$ and
$z_2=b\varepsilon\theta_2<\pi/4$.
The odd matching relation gives
\[
 -\theta_2\cot\theta_2
 =\varepsilon\theta_2\cot(b\varepsilon\theta_2)
 \le 1/b=-a\cot a.
\]
Since $s\mapsto-s\cot s$ is strictly increasing on $(\pi/2,\pi)$,
this proves $\theta_2\le a$. Define
\[
 d_1=\pi^2/4-\theta_1^2,
 \qquad d_2=a^2-\theta_2^2.
\]
Theorem~\ref{thm:lemAdp} and the deficit identity yield
$0\le d_2\le d_1$.
The even matching equation, with $\delta_1=\pi/2-\theta_1$, gives
\[
 0<\delta_1
 =\arctan\bigl(\varepsilon\tan(b\varepsilon\theta_1)\bigr)
 \le\varepsilon\tan(\tfrac\pi2 b\varepsilon).
\]
Consequently
\begin{equation}
 0\le G(R,u)-\bar D(u)
 =\frac{d_1-d_2}{u^2}
 \le\frac{\pi\varepsilon}{u^2}
       \tan(\tfrac\pi2 b\varepsilon)
 =O_u(R^{-1}).
 \label{eq:r8-fixed-u-bound}
\end{equation}
In particular $G(R,u^*)\to M$, so
$\limsup_{R\to\infty}Rm_R\le M$.

For every $R\ge1500$ and every $0<u<1/2$, either $w\le2$, when
Lemma~\ref{lem:sliver} gives $G>15\pi^2/4>3\pi^2>M$, or $w\ge2$,
when Theorem~\ref{thm:lemAdp} gives $G\ge\bar D(u)\ge M$.
Taking the infimum proves $Rm_R\ge M$ and hence the asserted limit.
Evaluating~\eqref{eq:r8-fixed-u-bound} at $u^*$ also proves
$0\le Rm_R-M=O(R^{-1})$.

For the stated near-minimizers,
\[
 R_jm_{R_j}\le G(R_j,u_j)
 \le R_jm_{R_j}+R_j\eta_j\longrightarrow M.
\]
The strict sliver margin implies $w_j=u_j\sqrt{R_j}>2$ eventually.
The large-$w$ lower bound then gives
\[
 M\le\bar D(u_j)\le G(R_j,u_j)\longrightarrow M.
\]
Any subsequential limit of $u_j$ in $[0,1/2]$ is interior:
a limit of zero contradicts the infinite left endpoint limit of
$\bar D$, while a limit of $1/2$ contradicts its right endpoint limit
$3\pi^2>M$. For an interior limit, continuity and the uniqueness in
Theorem~\ref{thm:T2} force that limit to be $u^*$.
Thus the entire sequence converges to $u^*$.
\end{proof}

% Integration note: remove old claims that script 16 certifies this region.
% Historical scripts remain immutable; their PASS labels are not used here.
% This fragment proves the stated symmetric three-block-family result.
% It does not extend the optimization to arbitrary densities or certify T3.
